% ── Chapter 2: Related Work and Background Theory ─────────────
\chapter{งานวิจัยที่เกี่ยวข้องและทฤษฎีพื้นฐาน (Related Work and Background Theory)}

\section{การทดสอบความไวต่อยาปฏิชีวนะ (Antibiotic Susceptibility Testing)}

การทดสอบ Kirby-Bauer disk diffusion assay ได้รับการอธิบายและกำหนดมาตรฐานอย่างเป็นระบบครั้งแรกโดย Bauer, Kirby, Sherris และ Turck ในปี 1966 และตั้งแต่นั้นมาก็ได้กลายเป็นเทคนิคที่ได้รับการยอมรับอย่างกว้างขวางที่สุดสำหรับการทดสอบความไวต่อยาต้านจุลชีพในห้องปฏิบัติการจุลชีววิทยาทางคลินิกทั่วโลก \cite{asm2009kirby,clsi2025} ความยั่งยืนของวิธีนี้เกิดจากต้นทุนอุปกรณ์ที่ต่ำ ความเรียบง่ายของขั้นตอน และความสามารถในการปรับใช้กับเชื้อก่อโรคแบคทีเรียที่เติบโตในสภาวะที่มีออกซิเจนได้เกือบทุกชนิด วิธีนี้ขึ้นอยู่กับหลักการของการแพร่กระจายของยาปฏิชีวนะผ่านวุ้นกึ่งแข็ง: เมื่อแผ่นดิสก์กระดาษที่ชุ่มด้วยยาปฏิชีวนะถูกวางลงบนพื้นผิวของจานวุ้นที่ปลูกเชื้อไว้ โมเลกุลของยาจะแพร่ออกไปในแนวรัศมีจากแผ่นดิสก์ สร้างระดับความเข้มข้นที่ลดลงแบบเอกซ์โพเนนเชียลตามระยะทาง ในจุดที่ความเข้มข้นของยาปฏิชีวนะในพื้นที่สูงกว่าค่าความเข้มข้นต่ำสุดที่ยับยั้งเชื้อได้ (minimum inhibitory concentration หรือ MIC) ของสิ่งมีชีวิตที่ทดสอบ การแพร่กระจายของแบคทีเรียจะถูกยับยั้ง ทำให้เกิดวงใสทรงกลมบนสนามการเจริญเติบโตของแบคทีเรียที่หนาแน่น

เส้นผ่านศูนย์กลางของโซนการยับยั้ง (zone of inhibition หรือ ZOI) นี้มีความสัมพันธ์แบบผกผันกับ MIC: สิ่งมีชีวิตที่มีความไวสูง (MIC ต่ำ) จะแสดงโซนขนาดใหญ่ ในขณะที่สิ่งมีชีวิตที่ดื้อยาโดยกำเนิดหรือดื้อยาที่ได้รับมาจะแสดงโซนขนาดเล็กหรือไม่มีเลย ความสัมพันธ์ระหว่างเส้นผ่านศูนย์กลาง ZOI และ MIC นี้ไม่ใช่แบบเส้นตรง แต่เป็นไปตามความสัมพันธ์ทางคณิตศาสตร์ที่ขึ้นอยู่กับความลึกของวุ้น ความแรงของแผ่นดิสก์ สัมประสิทธิ์การแพร่ และระยะเวลาการฟักตัว ซึ่งทั้งหมดนี้ได้รับการควบคุมอย่างระมัดระวังโดยโปรโตคอลมาตรฐาน สำหรับการใช้งานทางคลินิก การวัดเส้นผ่านศูนย์กลาง ZOI ที่ต่อเนื่องจะถูกแปลงเป็นการตีความความไวแบบแบ่งประเภท (Susceptible, Intermediate หรือ Resistant) โดยการใช้จุดตัดการตีความ (interpretive breakpoints) ที่ตีพิมพ์ ซึ่งตัวจุดตัดเองได้มาจากข้อมูลทางคลินิกและจุลชีววิทยา วิธีนี้ได้ถูกนำไปใช้กับแบคทีเรียที่มีความเกี่ยวข้องทางคลินิกในวงกว้าง รวมถึง \textit{Bacillus cereus} \cite{fitwi2020} และ \textit{Bacillus subtilis} \cite{jiang2023} ท่ามกลางเชื้อก่อโรคอื่นๆ อีกมากมายที่พบในการปฏิบัติงานในห้องปฏิบัติการตามปกติ

\subsection{มาตรฐานทางคลินิก: CLSI และ EUCAST (Clinical Standards: CLSI and EUCAST)}

จุดตัดการตีความ---เกณฑ์เส้นผ่านศูนย์กลางที่ใช้ในการจำแนกสิ่งมีชีวิตว่าเป็น ไวต่อยา (S), ระดับปานกลาง (I) หรือ ดื้อยา (R)---ได้รับการตีพิมพ์และปรับปรุงเป็นประจำโดยองค์กรกำหนดมาตรฐานระหว่างประเทศหลักสองแห่ง ได้แก่ Clinical and Laboratory Standards Institute (CLSI) ในสหรัฐอเมริกา และ European Committee on Antimicrobial Susceptibility Testing (EUCAST) ในยุโรป \cite{eucast2023} แม้ว่าทั้งสององค์กรจะมุ่งเป้าหมายทางคลินิกเดียวกัน แต่แนวทางระเบียบวิธีในการกำหนดจุดตัดนั้นแตกต่างกันในส่วนที่สำคัญ และเกณฑ์ที่ตีพิมพ์สำหรับคู่ยาและสิ่งมีชีวิตเดียวกันอาจแตกต่างกันได้หลายมิลลิเมตร

CLSI ตีพิมพ์มาตรฐาน disk diffusion ในเอกสาร M02 ซึ่งปัจจุบันอยู่ในฉบับที่สิบสี่ \cite{clsi2025} ในขณะที่มาตรฐานประสิทธิภาพการทดสอบความไวต่อยายังถูกกำหนดไว้ใน M100 ฉบับที่ 35 \cite{clsi_m100_2025} จุดตัดถูกกำหนดโดยกลุ่มทำงานสหวิทยาการที่รวบรวมข้อมูลการกระจาย MIC จากคอลเลกชันขนาดใหญ่ของเชื้อแยกทางคลินิก, การจำลอง PK/PD ของการได้รับยาในขนาดยาที่ใช้รักษามาตรฐาน และข้อมูลผลลัพธ์ทางคลินิกจากการทดลองแบบสุ่มที่มีกลุ่มควบคุม หมวดหมู่ระดับปานกลาง (I) ในระบบ CLSI ตามธรรมเนียมเดิมระบุถึงโซนกันชนที่ความไม่แน่นอนทางเทคนิคอาจนำไปสู่การจัดประเภทที่ผิดพลาด แต่รุ่นหลังๆ ได้ปรับตัวให้สอดคล้องกับแนวคิดของ EUCAST เกี่ยวกับหมวดหมู่ ``Susceptible, Increased Exposure'' (I) มากขึ้น ซึ่งระบุถึงความไวต่อยาเมื่อมีการรับยาในระดับสูงสุดผ่านการปรับปริมาณยาให้เหมาะสม

EUCAST ซึ่งตีพิมพ์ตารางจุดตัดเป็นประจำทุกปีพร้อมหมายเลขเวอร์ชัน (เวอร์ชันปัจจุบันที่อ้างอิงในโครงงานนี้คือ 13.1 \cite{eucast2023}) กำหนดจุดตัดจากเป้าหมาย pharmacokinetic/pharmacodynamic (PK/PD), ค่าจุดตัดทางระบาดวิทยา (epidemiological cutoff values หรือ ECOFFs) และข้อมูลผลลัพธ์ทางคลินิก \cite{giske2022} ค่า ECOFFs แสดงถึงขอบบนของการกระจาย MIC สำหรับประชากรแบคทีเรียสายพันธุ์ป่า (wild-type) ที่ไม่มีกลไกการดื้อยาที่ได้รับมา สิ่งมีชีวิตใดๆ ที่มี MIC สูงกว่า ECOFF จะถูกกำหนดให้เป็นสายพันธุ์ที่ไม่ใช่ป่า (non-wild-type หรือ NWT) และสันนิษฐานว่ามีตัวกำหนดการดื้อยา จุดตัดของ EUCAST มักจะรุนแรงกว่าในการแยกประชากรที่ไวต่อยาออกจากประชากรที่ดื้อยา ซึ่งสะท้อนถึงกรอบงานทางเภสัชวิทยาของยุโรปที่แตกต่างจากมาตรฐานของสหรัฐอเมริกา \cite{kahlmeter2019} จุดตัดทั้งสองชุดได้รับการดำเนินการในฐานข้อมูลของระบบปัจจุบัน ช่วยให้ห้องปฏิบัติการสามารถเลือกมาตรฐานที่เหมาะสมตามบริบทด้านกฎระเบียบและทางคลินิกของตน

\begin{figure}[H]
  \centering
  \safeimg[width=0.78\textwidth]{figures/_review/clsi_breakpoint_table.png}{CLSI breakpoint table for S. aureus}
  \caption{ตัวอย่างตารางเกณฑ์การแปลผลเส้นผ่านศูนย์กลางโซน (Zone diameter breakpoints) สำหรับ
    \textit{Staphylococcus aureus} ตามมาตรฐาน CLSI แสดงช่วง Resistant, Intermediate และ Susceptible
    สำหรับยาปฏิชีวนะแต่ละชนิด}
  \label{fig:clsi_table}
\end{figure}

\subsection{ความสำคัญทางคลินิกของความแม่นยำในการวัด (Clinical Significance of Measurement Precision)}

ผลที่ตามมาทางคลินิกของข้อผิดพลาดในการวัดในวิธี disk diffusion นั้นส่งผลโดยตรงและอาจรุนแรง เนื่องจากจุดตัดถูกแสดงเป็นเกณฑ์จำนวนเต็มในหน่วยมิลลิเมตร ข้อผิดพลาดในการวัดเพียงมิลลิเมตรเดียวสามารถเปลี่ยนการจำแนกจากไวต่อยาเป็นดื้อยาหรือในทางกลับกัน การจำแนกที่ผิดพลาดนี้มีสองรูปแบบที่อันตราย: ความไวต่อยาเท็จ (false susceptibility: การจัดประเภทสิ่งมีชีวิตที่ดื้อยาว่าไวต่อยา) นำไปสู่ความล้มเหลวในการรักษาและอาจส่งผลถึงชีวิต การดื้อยาเท็จ (false resistance: การจัดประเภทสิ่งมีชีวิตที่ไวต่อยาว่าดื้อยา) นำไปสู่การขยับไปใช้ยาในลำดับที่สองหรือสามโดยไม่จำเป็น ซึ่งส่งผลต่อแรงกดดันในการคัดเลือกเชื้อดื้อยาและต้นทุนที่สูงขึ้น Humphries et al.\ \cite{humphries2018} กำหนดค่าความคลาดเคลื่อนสมบูรณ์เฉลี่ยที่น้อยกว่า 3\,mm เป็นขอบเขตของความยอมรับได้ทางคลินิกสำหรับเครื่องมือวัด AST อัตโนมัติ และเกณฑ์นี้ทำหน้าที่เป็นเป้าหมายประสิทธิภาพหลักสำหรับโครงงานปัจจุบัน

\section{Deep Learning สำหรับการตรวจจับวัตถุ (Deep Learning for Object Detection)}

การตรวจจับวัตถุ (Object detection) ใน computer vision เป็นงานที่ดำเนินการระบุตำแหน่งและจำแนกวัตถุหนึ่งชิ้นหรือมากกว่าภายในภาพเดียวพร้อมกัน Machine learning ตามคำจำกัดความกว้างๆ รวมถึงกระบวนทัศน์การเรียนรู้แบบมีผู้สอน (supervised), แบบไม่มีผู้สอน (unsupervised) และการเรียนรู้แบบเสริมกำลัง (reinforcement learning) \cite{alam2023ml,murphy2025rl} ซึ่ง supervised deep learning ได้รับการพิสูจน์แล้วว่ามีประสิทธิภาพสูงสุดสำหรับงานการจดจำภาพ Convolutional Neural Networks (CNNs) ก่อให้เกิดกระดูกสันหลังของตัวตรวจจับวัตถุสมัยใหม่ส่วนใหญ่ \cite{purwono2023cnn} และการทบทวนแอปพลิเคชัน deep learning ทางการแพทย์อย่างเป็นระบบได้ยืนยันความเหมาะสมสำหรับงานการวิเคราะห์ภาพทางคลินิกที่หลากหลาย \cite{egger2022} ตัวตรวจจับต้องส่งออกข้อมูลสำหรับวัตถุแต่ละชิ้นที่ตรวจพบ ทั้ง bounding box ที่ระบุขอบเขตเชิงพื้นที่ของวัตถุและป้ายกำกับคลาสที่ระบุหมวดหมู่ นี่เป็นสิ่งที่ซับซ้อนกว่าการจำแนกรูปภาพ (image classification) พื้นฐานที่กำหนดป้ายกำกับเดียวให้กับรูปภาพทั้งหมด เนื่องจากจำนวนและตำแหน่งของวัตถุไม่ทราบล่วงหน้าและต้องสรุปจากเนื้อหาในภาพ แนวทาง Deep learning ได้ครอบงำเกณฑ์มาตรฐานการตรวจจับวัตถุนับตั้งแต่การตีความครอบครัว Region-based CNN (R-CNN) และโมเดล YOLO ซึ่งมีการปรับปรุงทั้งความแม่นยำและความเร็วในการประมวลผลอย่างต่อเนื่อง

ตัวตรวจจับวัตถุแบบ deep learning สมัยใหม่โดยทั่วไปจะแบ่งออกเป็นสองตระกูลสถาปัตยกรรม: ตัวตรวจจับแบบสองขั้นตอน (two-stage detectors) ซึ่งขั้นแรกจะสร้างข้อเสนอภูมิภาค (region proposals) แล้วจึงจำแนกข้อเสนอแต่ละรายการ และตัวตรวจจับแบบขั้นตอนเดียว (one-stage หรือ single-shot detectors) ซึ่งทำนาย bounding boxes และความน่าจะเป็นของคลาสโดยตรงในการส่งต่อข้อมูลผ่านเครือข่ายเพียงครั้งเดียวเหนือภาพเต็ม แต่ละตระกูลนำเสนอข้อดีข้อเสียที่แตกต่างกันระหว่างความแม่นยำในการตรวจจับ ต้นทุนการคำนวณ และความหน่วงในการประมวลผล

\subsection{Faster R-CNN}

Faster R-CNN (Region-based Convolutional Neural Network) ถูกนำเสนอโดย Ren et al.\ ในปี 2015 เป็นความก้าวหน้าครั้งสำคัญที่รวมการสร้างข้อเสนอภูมิภาคเข้ากับหัวการตรวจจับให้เป็นเครือข่ายเดียวที่สามารถฝึกฝนแบบครบวงจร (end-to-end) \cite{ren2015faster} ก่อนหน้า Faster R-CNN ไปป์ไลน์ของ R-CNN และ Fast R-CNN ต้องการขั้นตอนการค้นหาแบบเลือกเฟ้น (selective search) ที่แยกจากกันและมีต้นทุนการคำนวณสูงเพื่อสร้างข้อเสนอภูมิภาค ซึ่งเป็นคอขวดของความเร็วในการประมวลผล Faster R-CNN แทนที่สิ่งนี้ด้วย Region Proposal Network (RPN) ที่ใช้กระดูกสันหลังแบบคอนโวลูชันร่วมกับหัวการตรวจจับ ทำให้สามารถสร้างข้อเสนอได้เกือบเรียลไทม์

\begin{figure}[H]
  \centering
  \safeimg[width=0.82\textwidth]{figures/faster_rcnn_architecture.png}{Faster R-CNN two-stage detection architecture}
  \caption{สถาปัตยกรรมของ Faster R-CNN แสดง Convolutional backbone ที่ใช้ร่วมกัน,
    Region Proposal Network (RPN) ที่สร้างบริเวณวัตถุที่เป็นไปได้ และหัวการตรวจจับ
    ที่ทำการจำแนกและปรับแต่งข้อเสนอแต่ละรายการ}
  \label{fig:faster_rcnn}
\end{figure}

สถาปัตยกรรม Faster R-CNN ทำงานในสองขั้นตอน ในขั้นตอนแรก เครือข่ายคอนโวลูชันกระดูกสันหลัง (ในโครงงานนี้ใช้ ResNet-50 \cite{he2016deep}) จะสกัดแผนผังคุณลักษณะเชิงพื้นที่ที่เข้มข้นจากรูปภาพอินพุตแบบเต็ม จากนั้น RPN จะเลื่อนเครือข่ายขนาดเล็กไปบนแผนผังคุณลักษณะนี้ และในแต่ละตำแหน่งเชิงพื้นที่ จะทำนายคะแนนความเป็นวัตถุ (objectness scores) และค่าชดเชยการถดถอย bounding box สำหรับชุดของ anchor boxes ที่กำหนดไว้ล่วงหน้าซึ่งมีมาตราส่วนและอัตราส่วนภาพที่หลากหลาย Anchor boxes ที่มีคะแนนสูงกว่าเกณฑ์ความเป็นวัตถุจะถูกเก็บไว้เป็นข้อเสนอภูมิภาค (โดยทั่วไปคือ 256--512 ต่อภาพหลังจากทำ non-maximum suppression) ในขั้นตอนที่สอง การดำเนินการ RoI Pooling (หรือ RoI Align) จะสกัดเวกเตอร์คุณลักษณะขนาดคงที่สำหรับแต่ละข้อเสนอจากแผนผังคุณลักษณะกระดูกสันหลังที่แชร์ร่วมกัน และหัวการจำแนกประเภทที่เชื่อมต่ออย่างสมบูรณ์จะกำหนดป้ายกำกับคลาสและปรับแต่ง bounding box สำหรับแต่ละภูมิภาค

การใช้ ResNet-50 เป็นกระดูกสันหลังการสกัดคุณลักษณะมีความเหมาะสมอย่างยิ่งสำหรับโครงงานนี้ He et al.\ แสดงให้เห็นว่าการเชื่อมต่อแบบเรซิดูอัล (residual connections)---การเชื่อมต่อทางลัดที่ข้ามเลเยอร์คอนโวลูชันหนึ่งเลเยอร์หรือมากกว่า---ช่วยให้สามารถฝึกเครือข่ายที่ลึกขึ้นอย่างมีนัยสำคัญโดยไม่มีปัญหาการเสื่อมถอยของความชัน (gradient degradation) ที่เคยจำกัดความลึกไว้ \cite{he2016deep} ResNet-50 ประกอบด้วย 50 เลเยอร์ที่จัดระเบียบเป็นบล็อกเรซิดูอัล และการฝึกอบรมล่วงหน้า (pretraining) บน ImageNet ให้การเริ่มต้นของคุณลักษณะระดับต่ำ เช่น ขอบและพื้นผิว และคุณลักษณะรูปร่างระดับกลางที่ถ่ายโอนมายังขอบเขต AST ได้อย่างมีประสิทธิภาพ แม้ว่าจะมีรูปภาพ AST ที่ติดป้ายกำกับจำนวนค่อนข้างน้อยก็ตาม กระดูกสันหลัง ResNet-50 ในโครงงานนี้เริ่มต้นด้วยน้ำหนักที่ฝึกล่วงหน้าจาก ImageNet และปรับจูน (fine-tuned) แบบครบวงจรบนชุดข้อมูล AST

Faster R-CNN ได้รับการยอมรับในด้านความแม่นยำในการตรวจจับที่สูง โดยเฉพาะอย่างยิ่งในฉากที่ซับซ้อนซึ่งมีวัตถุที่ทับซ้อนกันหรืออยู่กันอย่างหนาแน่น---ซึ่งเป็นคุณลักษณะที่เกี่ยวข้องโดยตรงกับจาน AST ที่โซนจากแผ่นดิสก์ที่อยู่ติดกันอาจทับซ้อนกันบางส่วน สถาปัตยกรรมแบบสองขั้นตอนช่วยให้หัวการจำแนกประเภททำงานบนข้อเสนอที่มีคุณภาพสูงและได้รับการขัดเกลาเชิงพื้นที่ แทนที่จะเป็นตารางที่หนาแน่น ซึ่งช่วยลดอัตราการตรวจจับผิดพลาด (false positive) ข้อจำกัดหลักของ Faster R-CNN คือความเร็วในการประมวลผล: ไปป์ไลน์แบบสองขั้นตอนนำมาซึ่งความหน่วงที่ยอมรับได้สำหรับการวิเคราะห์แบบออฟไลน์ แต่อาจจำกัดปริมาณการประมวลผลแบบเรียลไทม์

\subsection{YOLOv11}

โมเดล YOLO (You Only Look Once) เป็นตัวตรวจจับแบบขั้นตอนเดียวที่กำหนดรูปแบบการตรวจจับวัตถุใหม่ให้เป็นปัญหาการถดถอย (regression) แบบรวมศูนย์ แทนที่จะสร้างข้อเสนอภูมิภาคและจำแนกแยกกัน โมเดล YOLO จะทำนายพิกัด bounding box และความน่าจะเป็นของคลาสโดยตรงจากรูปภาพเต็มในการส่งต่อข้อมูลผ่านเครือข่ายเพียงครั้งเดียว \cite{wang2024yolov11,khanam2024yolov11} การออกแบบนี้ช่วยขจัดขั้นตอนข้อเสนอภูมิภาคออกไปทั้งหมด ทำให้สามารถประมวลผลได้เร็วขึ้นอย่างมาก โดยแลกมาด้วยความแม่นยำในวัตถุขนาดเล็กหรือวัตถุที่ทับซ้อนกันอย่างหนาแน่นในเวอร์ชันแรกๆ การศึกษาเชิงเปรียบเทียบได้ยืนยันถึงประสิทธิภาพที่สามารถแข่งขันได้ของ YOLOv11 เมื่อเทียบกับรุ่นก่อนๆ และ Faster R-CNN ในงานการตรวจจับที่หลากหลาย \cite{sharma2024yolo_comparison}

\begin{figure}[H]
  \centering
  \safeimg[width=0.82\textwidth]{figures/yolov11_architecture.png}{YOLOv11 single-stage detection architecture}
  \caption{สถาปัตยกรรมของ YOLOv11 แสดงเครือข่าย Backbone สำหรับสกัดคุณลักษณะ,
    ส่วนคอ (neck) แบบ Feature Pyramid สำหรับการตรวจจับหลายมาตราส่วน
    และหัวการตรวจจับแยกส่วน (Decoupled heads) สำหรับการถดถอย Bounding box และการทำนายคลาส}
  \label{fig:yolov11}
\end{figure}

YOLOv11 ซึ่งเป็นเวอร์ชันล่าสุดในตระกูล YOLO ในขณะที่ทำโครงงานนี้และพัฒนาโดย Ultralytics ได้รวมความก้าวหน้าทางสถาปัตยกรรมหลายประการเหนือรุ่นก่อนหน้า \cite{wang2024yolov11} กระดูกสันหลังใช้บล็อก C3k2 ที่ปรับปรุงใหม่ (การออกแบบคอนโวลูชันที่ปรับพารามิเตอร์ใหม่) เพื่อเพิ่มประสิทธิภาพการสกัดคุณลักษณะ ส่วนคอ (neck) ใช้ Path Aggregation Feature Pyramid Network (PA-FPN) เพื่อหลอมรวมคุณลักษณะจากมาตราส่วนกระดูกสันหลังหลายระดับ ช่วยให้ตรวจจับวัตถุได้หลายขนาด ตั้งแต่แผ่นดิสก์ยาปฏิชีวนะขนาดเล็กไปจนถึงโซนการยับยั้งขนาดใหญ่ที่อาจครอบคลุมส่วนสำคัญของจาน หัวการตรวจจับถูกแยกออก (decoupled) โดยมีกิ่งก้านที่แยกจากกันสำหรับการทำนายการถดถอย bounding box และความน่าจะเป็นของคลาสอย่างเป็นอิสระ ซึ่งแสดงให้เห็นว่าช่วยปรับปรุงการลู่เข้า (convergence) และความแม่นยำในตัวตรวจจับแบบขั้นตอนเดียว

YOLOv11 มีให้เลือกห้าขนาด (n, s, m, l, x) ที่ให้ทางเลือกต่อเนื่องระหว่างความสามารถของโมเดลและความเร็วในการประมวลผล ในโครงงานนี้ได้เลือกใช้รุ่น nano (n) หรือ YOLOv11n เพื่อให้ความสำคัญกับความเร็วในการประมวลผลและความเหมาะสมสำหรับการติดตั้งใช้งานในอนาคตบนฮาร์ดแวร์ที่มีทรัพยากรจำกัด เช่น เวิร์กสเตชันห้องปฏิบัติการที่ไม่มี GPU เฉพาะ หรือฮาร์ดแวร์ปลายทาง (edge hardware) แม้จะเป็นรุ่นที่เล็กที่สุด แต่ YOLOv11n ก็บรรลุความแม่นยำในการตรวจจับที่แข่งขันได้ในชุดข้อมูล AST โดยได้รับประโยชน์จากการฝึกอบรมล่วงหน้าบนเกณฑ์มาตรฐาน COCO ขนาดใหญ่และการปรับจูนบนภาพ AST ที่ติดป้ายกำกับของโครงงาน โมเดลตรวจจับวัตถุสองคลาส ได้แก่ \texttt{antibiotic} (แผ่นดิสก์ที่ชุ่มด้วยยา ซึ่งใช้เป็นเป้าหมาย OCR และข้อมูลอ้างอิงการปรับเทียบ) และ \texttt{disk-zone} (โซนการยับยั้งโดยรอบ ซึ่งใช้สกัดการวัดเส้นผ่านศูนย์กลาง)

\subsection{Hough Circle Transform}

Hough Circle Transform เป็นเทคนิค computer vision แบบคลาสสิกสำหรับการตรวจจับโครงสร้างวงกลมในภาพโดยไม่ต้องใช้โมเดลที่เรียนรู้หรือข้อมูลการฝึกที่ติดป้ายกำกับ ทำงานโดยการแปลงภาพจากโดเมนเชิงพื้นที่ไปยังพื้นที่พารามิเตอร์ (parameter space) ที่กำหนดโดยสามมิติ: พิกัด $x$ และ $y$ ของจุดศูนย์กลางวงกลม และรัศมีวงกลม $r$ สำหรับจุดขอบแต่ละจุดที่ตรวจพบในภาพ (โดยทั่วไปผ่าน Canny edge detection) อัลกอริทึมจะลงคะแนน (vote) ให้กับวงกลมทั้งหมดในพื้นที่พารามิเตอร์ที่อาจสร้างจุดขอบนั้นได้ วงกลมที่สะสมคะแนนได้เหนือเกณฑ์ในตัวสะสมพื้นที่พารามิเตอร์จะสอดคล้องกับวงกลมจริงในภาพ แนวทางการตรวจจับวงกลมทางเลือกที่ใช้ learning automata \cite{cuevas2014}, deep learning \cite{ercan2020} และรุ่นที่ใช้การไล่ระดับสี (gradient-based) \cite{shehata2015hough} ได้ถูกนำเสนอเพื่อแก้ไขข้อจำกัดด้านความไวของทรานส์ฟอร์มแบบคลาสสิกภายใต้สภาวะภาพในโลกแห่งความเป็นจริง

\begin{figure}[H]
  \centering
  \safeimg[width=0.75\textwidth]{figures/hough_circle_diagram.png}{Hough Circle Transform parameter space voting diagram}
  \caption{ภาพประกอบ Hough Circle Transform จุดขอบแต่ละจุดจะลงคะแนนให้กับกรวยของ
    วงกลมที่เป็นไปได้ในพื้นที่พารามิเตอร์ $(x_c, y_c, r)$
    โดยวงกลมจริงจะสร้างยอดที่โดดเด่นในตัวสะสมคะแนน}
  \label{fig:hough}
\end{figure}

Hough Circle Transform มีประสิทธิภาพในการคำนวณและไม่สร้างผลการตรวจจับผิดพลาดสำหรับโซนที่เป็นวงกลมชัดเจนภายใต้สภาวะภาพที่ควบคุม อย่างไรก็ตาม มันมีความไวต่อแหล่งที่มาของการเสื่อมโทรมหลายอย่างที่พบบ่อยในภาพ AST จริง การสะท้อนของพื้นผิววุ้นที่ไม่สม่ำเสมอหรือแสงที่ไม่สม่ำเสมอสร้างการตอบสนองของขอบที่ผิดพลาดซึ่งสร้างสัญญาณรบกวนในตัวสะสม โซนที่ไม่ได้เป็นวงกลมที่สมบูรณ์---เนื่องจากการทับซ้อนของโซน ความไม่สม่ำเสมอของพื้นผิววุ้น หรือรูปแบบการเจริญเติบโตของแบคทีเรียที่ผิดปกติ---จะไม่สะสมคะแนนได้อย่างชัดเจน นำไปสู่การตรวจจับที่พลาดไปหรือการประมาณรัศมีที่ไม่แม่นยำ นอกจากนี้ ทรานส์ฟอร์มนี้ต้องการการปรับแต่งพารามิเตอร์หลายตัวด้วยมืออย่างระมัดระวัง (เช่น อัตราส่วนความละเอียดผกผัน \texttt{dp}, ระยะห่างขั้นต่ำระหว่างศูนย์กลาง \texttt{minDist} และเกณฑ์ขอบ Canny \texttt{param1} และ \texttt{param2}) ซึ่งมีปฏิสัมพันธ์กันในทางที่ซับซ้อน และค่าที่เหมาะสมที่สุดจะแตกต่างกันไปตามภาพที่มีมาตราส่วน แสง และชนิดของแบคทีเรียที่แตกต่างกัน ในโครงงานนี้ พารามิเตอร์ Hough Circle Transform ได้รับการปรับแต่งผ่านการค้นหาแบบตาราง (grid search) บนชุดข้อมูลเต็มรูปแบบ ซึ่งแสดงถึงสถานการณ์ที่ดีที่สุดสำหรับวิธีนี้ วิธีนี้ทำหน้าที่เป็นเกณฑ์มาตรฐานแบบคลาสสิกที่มีหลักการซึ่งโมเดลที่เรียนรู้จะถูกนำมาประเมินเปรียบเทียบ

\section{การจดจำอักขระด้วยแสงสำหรับฉลากแผ่นดิสก์ยาปฏิชีวนะ (Optical Character Recognition for Antibiotic Disk Labeling)}

แผ่นดิสก์ยาปฏิชีวนะจะมีรหัสตัวอักษรและตัวเลขสั้นๆ พิมพ์อยู่บนพื้นผิว ระบุชื่อยาปฏิชีวนะและความเข้มข้น (ตัวอย่างเช่น ``OX 1'' สำหรับ oxacillin 1\,$\mu$g หรือ ``AMX 25'' สำหรับ amoxicillin 25\,$\mu$g) การจดจำอักขระด้วยแสง (Optical Character Recognition หรือ OCR) ซึ่งประวัติการพัฒนาครอบคลุมความก้าวหน้าทางวิศวกรรมหลายทศวรรษ \cite{wang2023ocr_history} มอบวิธีการคำนวณเพื่ออ่านฉลากเหล่านี้โดยอัตโนมัติจากภาพถ่าย การอ่านฉลากเหล่านี้อย่างถูกต้องเป็นสิ่งสำคัญสำหรับการเชื่อมโยงแต่ละ ZOI ที่ตรวจพบกับยาปฏิชีวนะที่ถูกต้อง ซึ่งจะช่วยให้สามารถดึงจุดตัดที่ถูกต้องจากฐานข้อมูล CLSI/EUCAST ได้ นี่เป็นงาน OCR ที่ท้าทายด้วยเหตุผลหลายประการ: ข้อความที่พิมพ์มีขนาดเล็ก (โดยทั่วไปสูง 2--4 มม.), ความคมชัดระหว่างหมึกและพื้นผิวแผ่นดิสก์มีความแปรปรวนตามยี่ห้อของแผ่นดิสก์, พื้นผิวแผ่นดิสก์มีความโค้งเล็กน้อย, ภาพอาจถูกถ่ายในมุมที่ไม่เหมาะสม และแผ่นดิสก์หลายแผ่นอาจถูกบดบังบางส่วน

\subsection{EasyOCR และสถาปัตยกรรม CRAFT (EasyOCR and the CRAFT Architecture)}

EasyOCR \cite{jaidedai2024easyocr} เป็นเครื่องยนต์ OCR แบบเปิดเผยรหัสที่ใช้ deep learning ซึ่งรองรับภาษาและสคริปต์มากกว่า 80 ภาษา ขั้นตอนการตรวจจับข้อความขึ้นอยู่กับโมเดล CRAFT (Character Region Awareness for Text Detection) ที่นำเสนอโดย Baek et al.\ \cite{baek2019craft} CRAFT ตรวจจับภูมิภาคของอักขระแต่ละตัวและความสัมพันธ์ระหว่างอักขระที่อยู่ติดกัน แทนที่จะถือว่า bounding boxes ในระดับคำเป็นหน่วยการตรวจจับพื้นฐาน แนวทางการตรวจจับระดับอักขระนี้เหมาะสมอย่างยิ่งกับบริบทของฉลากยาปฏิชีวนะ ซึ่งต้องระบุตำแหน่งของสตริงตัวอักษรและตัวเลขสั้นๆ อย่างแม่นยำบนพื้นผิวแผ่นดิสก์ขนาดเล็กที่มีพื้นผิว

เครือข่าย CRAFT สร้างแผนผังเอาต์พุตสองแบบ: แผนผังคะแนนภูมิภาค (region score map: ระบุความน่าจะเป็นที่แต่ละพิกเซลเป็นศูนย์กลางของอักขระ) และแผนผังคะแนนความสัมพันธ์ (affinity score map: ระบุความน่าจะเป็นที่อักขระสองตัวที่อยู่ติดกันเป็นของคำเดียวกัน) แผนผังทั้งสองนี้จะถูกรวมเข้าด้วยกันเพื่อแบ่งส่วนอักขระแต่ละตัวและรวมกลุ่มกันเป็นอินสแตนซ์ข้อความระดับคำ ภูมิภาคอักขระที่ตรวจพบจะถูกส่งไปยังเครือข่ายการจดจำลำดับ---Convolutional Recurrent Neural Network (CRNN)---ซึ่งจะถ่ายทอดลำดับอักขระเป็นข้อความโดยใช้ตัวถอดรหัส connectionist temporal classification (CTC)

EasyOCR \cite{jaidedai2024easyocr} ได้รับการเปรียบเทียบกับ TyphoonOCR \cite{typhoonocr2025} ซึ่งเป็นโมเดล OCR ที่เน้นภาษาไทยที่พัฒนาโดย OpenTyphoon.ai สำหรับงานอ่านฉลากยาปฏิชีวนะ แม้ว่า TyphoonOCR จะให้ประสิทธิภาพที่แข็งแกร่งสำหรับข้อความสคริปต์ไทย แต่ประสิทธิภาพในสตริงตัวอักษรและตัวเลขละตินภายใต้เงื่อนไขที่ท้าทายของการถ่ายภาพแผ่นดิสก์นั้นด้อยกว่า EasyOCR \cite{baek2019craft,jaidedai2024easyocr} ดังนั้น EasyOCR จึงถูกเลือกเป็นเครื่องยนต์ OCR สำหรับไปป์ไลน์ขั้นสุดท้าย

\section{เทคนิคการเตรียมการประมวลผลภาพ (Image Preprocessing Techniques)}

การเตรียมการประมวลผลภาพเป็นตัวขับเคลื่อนที่สำคัญของทั้งโมเดลการตรวจจับวัตถุและไปป์ไลน์ OCR ซึ่งส่งผลโดยตรงต่ออัตราส่วนสัญญาณต่อสัญญาณรบกวนของคุณลักษณะที่สกัดโดยอัลกอริทึมปลายน้ำ \cite{qi2021image_enhancement} เทคนิคการแบ่งส่วนภาพ \cite{jasim2021segmentation,dhanachandra2015} และวิธีการแปลงเป็นภาพขาวดำ (binarization) ของเอกสาร \cite{yang2024binarization} ให้ข้อมูลสำหรับการเลือกการเตรียมการประมวลผลในโครงงานนี้ เพื่อจัดการกับการเสื่อมโทรมของภาพเฉพาะที่พบในการถ่ายภาพ AST ในห้องปฏิบัติการ

\begin{itemize}
  \item \textbf{Global thresholding (วิธีของ Otsu):} วิธีของ Otsu จะกำหนดเกณฑ์ความเข้มที่เหมาะสมที่สุดโดยอัตโนมัติสำหรับการแปลงรูปภาพโทนสีเทาเป็นภาพขาวดำ (binary) โดยการลดความแปรปรวนภายในคลาสระหว่างพิกเซลเบื้องหน้า (ข้อความ) และเบื้องหลัง วิธีนี้ใช้กับภาพที่ครอบเฉพาะแผ่นดิสก์ก่อนทำ OCR เพื่อแยกอักขระหมึกออกจากพื้นหลังพื้นผิวแผ่นดิสก์ ซึ่งช่วยปรับปรุงคุณภาพการแบ่งส่วนอักขระ

  \item \textbf{Contrast Limited Adaptive Histogram Equalization (CLAHE):} แตกต่างจากการปรับฮิสโตแกรมให้เท่ากันแบบโกลบอลที่คำนวณฟังก์ชันการจับคู่เดียวสำหรับทั้งภาพ CLAHE จะใช้การปรับฮิสโตแกรมให้เท่ากันอย่างอิสระกับกระเบื้อง (tiles) ขนาดเล็กของภาพ (พร้อมขีดจำกัดความคมชัดและขนาดกระเบื้องที่กำหนดค่าได้) และประมาณค่าในช่วงระหว่างกระเบื้องอย่างราบรื่น สิ่งนี้ช่วยปรับความคมชัดในระดับพื้นที่ให้สม่ำเสมอในภูมิภาคของภาพที่มีความสว่างเฉลี่ยต่างกันมาก---ซึ่งเป็นเหตุการณ์ทั่วไปในภาพถ่ายห้องปฏิบัติการที่จุดศูนย์กลางของจานวุ้นสว่างกว่าส่วนขอบเนื่องจากรูปทรงของแสงเหนือศีรษะ CLAHE ถูกนำไปใช้กับภาพแผ่นดิสก์ที่ครอบมาเป็นขั้นตอนการเตรียมการประมวลผลก่อน OCR เพื่อปรับปรุงทัศนวิสัยของข้อความที่มีความคมชัดต่ำ

  \item \textbf{Gaussian Blur:} ตัวกรองความถี่ต่ำ Gaussian จะถูกนำไปใช้กับรูปภาพจานแบบเต็มก่อนการตรวจจับขอบ (สำหรับ Hough Transform) และก่อนการหาเกณฑ์ ช่วยลดสัญญาณรบกวนเซ็นเซอร์ความถี่สูงที่อาจสร้างการตอบสนองขอบที่ผิดพลาด ขนาดของเคอร์เนลการเบลอจะถูกเลือกเพื่อรักษามาตราส่วนที่เกี่ยวข้องของขอบเขตโซนในขณะที่ระงับสัญญาณรบกวน

  \item \textbf{การดำเนินการทางสัณฐานวิทยา (dilation และ erosion):} การขยาย (dilation) ทางสัณฐานวิทยาจะขยายภูมิภาคเบื้องหน้า (ข้อความ) ในขณะที่การกัดกร่อน (erosion) จะทำให้เล็กลง เมื่อนำมาใช้ตามลำดับ (เป็นการดำเนินการเปิดหรือปิด) สิ่งเหล่านี้สามารถเติมช่องว่างเล็กๆ ในอักขระหรือลบจุดรบกวนขนาดเล็กจากภาพขาวดำได้ ในไปป์ไลน์การเตรียมการประมวลผล OCR การเลือกและขอบเขตของการดำเนินการทางสัณฐานวิทยาจะถูกนำไปใช้อย่างเหมาะสมตามอัตราส่วนพิกเซลข้อความที่ประมาณไว้ในภาพที่ครอบมา: ภาพที่มีพิกเซลข้อความเบาบางมาก (อักขระบาง) จะถูกขยายเพื่อกู้คืนการเชื่อมต่อของเส้น ในขณะที่ภาพที่มีพิกเซลหนาแน่นเกินไป (เส้นหนาหรืออักขระกลืนกัน) จะถูกกัดกร่อน

  \item \textbf{การแก้ไขการหมุนของภาพ (Image rotation correction):} ทิศทางของแผ่นดิสก์ยาปฏิชีวนะบนจานไม่ได้ถูกควบคุมระหว่างการตั้งค่าในห้องปฏิบัติการ---แผ่นดิสก์อาจถูกวางในมุมใดก็ได้ ส่งผลให้ฉลากที่พิมพ์บนแผ่นดิสก์อาจปรากฏหมุนที่มุมใดก็ได้เมื่อเทียบกับเฟรมภาพ เพื่อจัดการสิ่งนี้ ไปป์ไลน์ OCR จะใช้ EasyOCR ที่สี่ทิศทางมาตรฐาน (0$^\circ$, 90$^\circ$, 180$^\circ$ และ 270$^\circ$) และเลือกผลลัพธ์การจดจำที่มีความเชื่อมั่นสูงสุดจากทั้งสี่ทิศทาง
\end{itemize}

\section{งานวิจัยที่เกี่ยวข้องในการวิเคราะห์ AST อัตโนมัติ (Related Work in Automated AST Analysis)}

\subsection{แนวทาง Deep Learning สำหรับการวิเคราะห์ภาพ AST (Deep Learning Approaches for AST Image Analysis)}

การประยุกต์ใช้ deep learning และ computer vision ในการวิเคราะห์ภาพ AST เป็นสาขาการวิจัยที่มีความเคลื่อนไหวในช่วงทศวรรษที่ผ่านมา โดยได้รับแรงผลักดันจากการผสมผสานระหว่างความต้องการทางคลินิกและความก้าวหน้าของอัลกอริทึมที่เป็นแรงจูงใจให้เกิดโครงงานนี้ บทบาทที่กว้างขึ้นของ AI และ machine learning ในการทำนายและต่อสู้กับการดื้อยาต้านจุลชีพได้รับการทบทวนอย่างครอบคลุมโดย Bilal et al.\ \cite{bilal2025} ในขณะที่ Liao et al.\ \cite{liao2025} ได้สำรวจความก้าวหน้าโดยเฉพาะในงานวิจัยการทดสอบความไวต่อยาที่ขับเคลื่อนด้วย AI
Pascucci et al.\ \cite{pascucci2021,pascucci2020preprint} ได้พัฒนาแอปพลิเคชันมือถือบนฐาน AI ตัวแรกสำหรับการทดสอบการดื้อยาปฏิชีวนะอัตโนมัติ แสดงให้เห็นว่ากล้องสมาร์ทโฟนรวมกับ deep learning สามารถสร้างการวัด ZOI ที่ยอมรับได้ทางคลินิกในสภาพแวดล้อมที่มีทรัพยากรจำกัด

\begin{figure}[H]
  \centering
  \safeimg[width=\textwidth]{figures/_review/related_work_existing_system.png}{Existing automated AST system}
  \caption{ตัวอย่างระบบวิเคราะห์ AST อัตโนมัติเชิงพาณิชย์ที่มีอยู่แล้ว แสดงผลการตรวจจับ
    วงยับยั้งพร้อมตารางผลการจำแนก S/I/R ที่ใช้เป็นเกณฑ์มาตรฐานในการเปรียบเทียบ}
  \label{fig:related_existing}
\end{figure}
Rajasekar et al.\ \cite{rajasekar2022} ได้นำเสนอ Ab.ai ซึ่งเป็นเครื่องมือ AI เฉพาะทางสำหรับการรายงาน antibiograms โดยอัตโนมัติ บรรลุความแม่นยำสูงบนจาน disk diffusion ในสภาวะห้องปฏิบัติการมาตรฐาน
Yu et al.\ \cite{yu2025} แสดงให้เห็นถึงวิธี hue-contrast ที่ขับเคลื่อนด้วย AI สำหรับการระบุโซนการยับยั้ง ซึ่งเป็นทางเลือกนอกเหนือจากการถดถอย bounding-box
Callebaut et al.\ \cite{callebaut2024} ประเมินระบบ Radian\textregistered{} in-line carousel สำหรับการทดสอบ disk diffusion อัตโนมัติ กำหนดเกณฑ์มาตรฐานประสิทธิภาพเทียบกับการอ่านด้วยตนเองโดยผู้เชี่ยวชาญซึ่งโครงงานนี้นำมาใช้เป็นเป้าหมายเปรียบเทียบ

Rajasekar et al.\ \cite{rajasekar2023} ได้ดำเนินการทบทวนอย่างเป็นระบบเกี่ยวกับวิธีการวิเคราะห์ภาพบนฐาน AI ที่ตีพิมพ์สำหรับ AST และพบว่าโมเดล deep learning มีประสิทธิภาพเหนือกว่าแนวทางการประมวลผลภาพแบบคลาสสิกอย่างสม่ำเสมอและมีนัยสำคัญ---รวมถึง Hough Transform, active contours และ template matching---โดยเฉพาะอย่างยิ่งในสถานการณ์ที่มีความเกี่ยวข้องมากที่สุดกับการปฏิบัติทางคลินิก: จานที่มีแผ่นดิสก์จำนวนมาก, โซนที่ทับซ้อนกัน และการเจริญเติบโตของแบคทีเรียที่ไม่เป็นเนื้อเดียวกัน การทบทวนได้สำรวจวิธีการตั้งแต่อดีตที่มีการใช้เครือข่ายคอนโวลูชันดั้งเดิมกับภูมิภาคแผ่นดิสก์ที่ครอบมา ไปจนถึงไปป์ไลน์การตรวจจับแบบครบวงจรที่คล้ายกับที่พัฒนาขึ้นในโครงงานนี้

ท่ามกลางสิ่งที่ค้นพบที่สำคัญของ Rajasekar et al.\ คือการเพิ่มประสิทธิภาพที่มากที่สุดจาก deep learning เกิดขึ้นโดยเฉพาะในการจัดการกับโซนที่ทับซ้อนกัน---ซึ่งเป็นสถานการณ์ที่ทำลายข้อสมมติฐานของรูปร่างวงกลมที่แยกจากกันซึ่งเป็นพื้นฐานของแนวทาง Hough Transform---และความทนทานต่อความแปรปรวนของภาพถ่าย เช่น แสงที่ไม่สม่ำเสมอ ความผิดเพี้ยนของมุมมอง และสิ่งแปลกปลอมจากการบีบอัดภาพ \cite{rajasekar2023} สิ่งเหล่านี้ช่วยยืนยันการเลือกทางสถาปัตยกรรมของโครงงานปัจจุบัน: การใช้การตรวจจับวัตถุด้วย deep learning เป็นแนวทางหลัก และจัดให้ Hough Transform เป็นเพียงเกณฑ์มาตรฐานแบบคลาสสิกมากกว่าที่จะเป็นส่วนประกอบหลักของระบบ การทบทวนยังระบุด้วยว่าขนาดที่ค่อนข้างเล็กของชุดข้อมูล AST ส่วนใหญ่ที่ตีพิมพ์ (โดยทั่วไปต่ำกว่า 500 ภาพ) เป็นปัจจัยจำกัดประสิทธิภาพของ deep learning ซึ่งเป็นแรงจูงใจให้โครงงานปัจจุบันลงทุนในการสร้างข้อมูลสังเคราะห์เพื่อให้ถึงชุดข้อมูลฝึกอบรมมากกว่า 2,000 ภาพ

\subsection{มาตรฐานการประเมินสำหรับวิธี AST อัตโนมัติ (Evaluation Standards for Automated AST Methods)}

Humphries et al.\ \cite{humphries2018} ได้กล่าวถึงคำถามเชิงระเบียบวิธีเกี่ยวกับวิธีการประเมินระบบการวัด AST อัตโนมัติอย่างเหมาะสมเมื่อเทียบกับมาตรฐานทองคำ (gold standard) ของการอ่านด้วยตนเองโดยนักจุลชีววิทยาที่ผ่านการฝึกอบรม เอกสารแนวทางปฏิบัติที่ดีที่สุดของพวกเขา ซึ่งตีพิมพ์โดย CLSI Methods Development and Standardization Working Group ได้กำหนดเกณฑ์การประเมินที่สำคัญหลายประการที่โครงงานปัจจุบันนำมาใช้ สิ่งที่สำคัญที่สุดสำหรับงานนี้คือพวกเขาได้กำหนดว่าค่าความคลาดเคลื่อนสมบูรณ์เฉลี่ย (MAE) ที่น้อยกว่า 3\,mm ระหว่างการวัดเส้นผ่านศูนย์กลาง ZOI แบบอัตโนมัติและแบบแมนนวล แสดงถึงเกณฑ์ของความยอมรับได้ทางคลินิก---กล่าวคือ ต่ำกว่าระดับข้อผิดพลาดนี้ ความแตกต่างในการวัดเชิงระบบไม่น่าจะนำไปสู่การจำแนก S/I/R ที่แตกต่างกันสำหรับคู่จุดตัดของยาและสิ่งมีชีวิตส่วนใหญ่ที่ใช้อยู่ในปัจจุบัน

Humphries et al.\ ยังเน้นย้ำถึงความสำคัญของการประเมินวิธีการอัตโนมัติในคอลเลกชันที่หลากหลายของสายพันธุ์แบคทีเรีย ชั้นของยาปฏิชีวนะ และการกำหนดค่าจาน เนื่องจากข้อผิดพลาดที่ไม่มีนัยสำคัญสำหรับคู่ยาและสิ่งมีชีวิตหนึ่งอาจตัดสินผลทางคลินิกสำหรับอีกคู่หนึ่งได้ \cite{humphries2018} คำแนะนำนี้ให้ข้อมูลแก่การออกแบบการประเมินของโครงงานปัจจุบัน ซึ่งรายงาน MAE ในชุดข้อมูลทดสอบที่แยกไว้ทั้งหมด 235 ภาพที่นำมาจากคอลเลกชันจานจริง นอกจากนี้ Humphries et al.\ ยังระบุว่าความลำเอียงเชิงระบบ (ทิศทาง) มีความสำคัญพอๆ กับการระบุข้อผิดพลาดแบบสุ่ม เนื่องจากการประเมินเส้นผ่านศูนย์กลางโซนที่สูงหรือต่ำเกินไปอย่างสม่ำเสมอสามารถแก้ไขได้ด้วยขั้นตอนการแก้ไขความลำเอียงหลังการประมวลผล---ซึ่งเป็นสิ่งที่ขั้นตอนการปรับแต่งการปรับเทียบของระบบปัจจุบันดำเนินการอย่างแม่นยำ
